\documentclass[12pt]{article}
\usepackage[utf8]{inputenc}
\usepackage[T1]{fontenc}
\usepackage{hyperref}
\title{Papuma Language Resources for K-12 Curriculum Development}
\author{Compiled by Claude Sonnet 5.5 (claude-sonnet-5-5) for Andr\'as Kornai}
\date{}
\begin{document}
\maketitle

\section{Introduction}
Papuma (ISO 639-3 ppm) is an Austronesian language of the Yapen group, spoken in Papuma village on the south coast of Yapen Island, Papua, Indonesia, by about 600 people. Very little beyond encyclopedic entries exists; this report is correspondingly short.

\section{Primary Educational Resources}
None found.

\section{Linguistic and Reference Materials}
\begin{itemize}
\item Wikipedia, Papuma language: \url{https://en.wikipedia.org/wiki/Papuma_language}
\item Numeral word list entry (Max Planck EVA, Chan numerals database): \url{https://lingweb.eva.mpg.de/channumerals/Papuma.htm}
\end{itemize}

\section{Religious and Community Resources}
\begin{itemize}
\item Joshua Project people-group profile (note: it describes the language as Papuan, contradicting the Austronesian classification): \url{https://joshuaproject.net/people_groups/14310}
\end{itemize}

\section{Digital and Audio Resources}
None found.

\section{Curriculum Development Context}
No school materials, orthography or documentation beyond word lists were found. Curriculum work would require new documentation with the community.
\end{document}